\section{Further research questions and proposed programme}
\label{sec:agenda}

The following questions are concrete continuations of the proved
results. Each item identifies the missing mathematical step, so that a
future finite experiment can be judged against a defined goal.

\subsection{Angular geometry and fractional kernels}

\paragraph{1. Complete the original global normalized-radius conjecture.}
Prove or refute that $\eta_{a,b}(\rho)$ decreases on $(0,1]$ for
every integer $a\ge2,b>0$, and establish the proposed $a=1$
trichotomy. Theorem~\ref{thm:radius-local-sign} removes the first local
obstruction. A useful next step is an exact sign formula for
$\partial_\rho\eta$ evaluated at its implicit zero, followed by a
variation-diminishing or ordered-kernel argument. A proof must control
the moving zero, not merely the integrand at a fixed angle.

\paragraph{2. Classify fractional radial turning points.}
For $0<b<1$, the curve $a_*(b)$ determines the local quadratic
coefficient. Theorems~\ref{turn:thm:maximum} and \ref{turn:thm:minimum}
prove opposite local bifurcations on this curve, so the quartic sign
is not constant. Determine the number of zeros of
$Q(a_*(b),b)$, certify the diagnostic crossing near
$b=0.9974937898734204$, and analyze the sixth-order coefficient there.
This is a specific candidate for a higher-order radial transition.
Determine also whether a single global radial curve can have several
turning points; the local theorems do not exclude that possibility.

\paragraph{3. Determine the complex zero region below the moment boundary.}
The finite signed-measure classification is sharp, but it does not
classify all parameters for complex nonvanishing. For $a=0$,
$F_{0,b}(z)=z\Li_b(z)/(1-z)$ gives a different route. Determine
whether nonzero complex zeros appear for $0<a<1-b$, and, if they do,
locate the first parameter boundary. A generalized Stieltjes transform
or a distributional moment representation may replace finite variation;
it would require a new zero argument.

\paragraph{4. Resolve the two-parameter corner $a\downarrow0,b\to1$.}
For fixed $b\ne1$ the crossing approaches one by a power law, while at
$b=1$ its distance is essentially $e^{-1/a}$. Derive a uniform
interpolation when $b-1$ and $a$ vanish together. The exact rescaled
kernel is a useful starting point, but neither fixed-$b$ asymptotic
is uniform at this corner. Determine the corresponding total-variation
profile and the weak limits after renormalization.

\paragraph{5. Optimize the Euler constant in the extended domain.}
The constants $(1-b)^{-1}$, $\zeta(b)$ and
$e^{a(\gamma+\psi(a))}/(a^2\Gamma(a))$ give explicit enclosures.
They are not claimed to be optimal remainder constants. Characterize
the exact signed mass as a function of $(a,b)$, sharpen the remainder
using its vanishing moments, and design an interval evaluator that
remains effective as the positive density collapses into an atom.

\subsection{Lerch zeros and overlapping asymptotic regimes}

\paragraph{6. Unfold the central Bessel zero.}
For fixed $r>0$, $\Psi_r$ has a multiple zero at zero. Its splitting
depends on a spectral contribution that is invisible to every algebraic
order in $1/k$. Obtain its leading exponential scale and sign, including
possible cancellation among shifted terms. This should connect the
fixed-$(n,k)$ small-$\rho$ Puiseux theorem with the joint limit
$n=k+r$, but the limits need not commute.

\paragraph{7. Allow the Bessel order or zero label to grow.}
Determine a range of $r=r(k)$ and $m=m(k)$ in which the fixed-order
Bessel approximation is still uniform. Find the transition to an
appropriate large-order or turning-point profile. The first technical
obligation is a uniform coefficient bound for the symmetric product
after the factorial transform, rather than a fixed-$r$ formal expansion.

\paragraph{8. Sharpen the exponential Lerch deformation.}
The factor $e^{-k/20}$ is explicit and adequate for separation of
scales, but it is deliberately conservative. Optimize the coefficient
contour and derive the actual leading exponential term, including its
oscillatory sign. This would explain the changes of sign seen in the
full-Lerch diagnostics and quantify when the first shift dominates.

\paragraph{9. Determine the sign pattern of $f_{n,k}(2)$.}
The all-pair small-$\rho$ theorem reduces the remaining count choice
to a nonzero integer polynomial evaluated at $\log2$. Find a useful
closed sign criterion in large regions of the $(n,k)$ lattice, or an
asymptotic sign law when both indices grow. Transcendence proves
nonvanishing, but supplies no practical separation bound. Rational
interval certificates can produce a reliable initial map.

\subsection{Reflected moments and uniform asymptotics}

\paragraph{10. Extend the critical transition to complete finite orders.}
The gamma localization argument and exact moments give an algorithm
for higher terms. Compute and prove a useful recurrence for the
coefficient functions beyond $P_1$, determine whether all negative
powers of $L$ cancel, and produce explicit computable remainder
constants on prescribed compact $c$ intervals. The first cancellation
is proved here; a general cancellation pattern needs its own argument.

\paragraph{11. Treat proportional growth and symmetry of the saddle.}
For $m/n\to\lambda>0$, analyze the critical points of
\[
 \log\log\Gamma(x)+\lambda\log\log\Gamma(1-x).
\]
Prove uniqueness or classify competing maxima, obtain a uniform
Laplace expansion, and connect it to the $m\asymp n/\log n$
boundary regime. The exact symmetry $M_{n,m}=M_{m,n}$ is an
essential consistency test. Numerical maximization alone cannot
exclude an additional saddle.

\paragraph{12. Make optimal truncation uniform in the reflected exponent.}
The source's sharp fixed-$m$ remainder is a separate theorem from the
present joint transition. Determine a common range in $(n,m,N)$ for
which late inverse-function coefficients and the least-term remainder
are uniform. The early-pole resummation suggests the correct matching,
but does not license summation of a divergent infinite pole series.

\subsection{Formal relation quotients and special values}

\paragraph{13. Prove the frozen $S_6$ depth-two formula.}
Construct simultaneous Cayley and conjugation eigenvector generators,
reduce the weight-seven problem to the 7518-dimensional imaginary
Cayley quotient, and impose double shuffle there with exact lifting
data. A successful output is a rational certificate whose expanded
residual is exactly zero. A negative answer for a specified presentation
requires a separating functional; it does not disprove the numerical
identity unless the presentation is known to be complete.

\paragraph{14. Determine the intersection with double shuffle.}
Theorem~\ref{cayley:thm:ideal} answers the Cayley-only problem in all
weights. Find the Hilbert or character series after adding a precisely
specified regularized double-shuffle ideal. Compare raw rows,
multiplicative closures, distribution, and rational pullback relations.
Any dimension conjecture should specify which algebra and which ideal
are meant, and should distinguish formal dimensions from periods.

\paragraph{15. Generalize the involution quotient theorem.}
The proof uses a graded polynomial shuffle algebra, two commuting
involutions, and a finite endpoint correction. Extend it to other
root-of-unity alphabets and rational changes of variables when the
admissible subalgebra is stable. Determine which transformations can be
simultaneously linearized on generators and what replaces the dyadic
functional equation when the acting symmetry has higher order.

\paragraph{16. Search for short consequences of the explicit $S_p$ family.}
The all-depth expression is exact but expands exponentially. Seek
families in which conjugation, parity, or shuffle symmetries collapse
most terms, and identify weights where a depth-two basket can plausibly
be complete. Freeze each candidate basket and coefficient vector before
verification, retain exact residual enclosures, and then seek symbolic
proofs. The rejected earlier $S_8$ vector should remain a regression
case rather than being silently replaced.

\subsection{Suggested order of work}
The most immediate finite task is the equivariant Cayley reduction for
$S_6$, because the quotient theorem gives a concrete smaller target.
On the analytic side, the global integer-order radial conjecture and
the two-parameter kernel corner are natural next goals: the present
sign and scaling formulas expose their remaining obstacles. The central
Bessel unfolding and uniform late-term gamma theory are more demanding
matching problems. All four directions benefit from keeping exact
finite verification separate from conjectural global extrapolation.
